%! Author = sriadityavedantam
%! Date = 3/25/26

\section{Homework 3}

\begin{problem}[Infinite Products]{
    Consider
    \[
        \prod_{n=1}^\infty(1+a_n)
        \qquad\text{where } a_n\geq0.
    \]

    \begin{enumerate}[label={(\alph*)}]
\item Find an explicit formula for the partial products when $a_n=1/n$ and decide whether they converge.
Also examine the case $a_n=1/n^2$.
\item Show that the partial products converge if and only if $\sum a_n$ converges.
    \end{enumerate}
}
    \textbf{(a).}
    If $a_n=\frac1n$, then
    \[
        p_m=\prod_{n=1}^m\left(1+\frac1n\right)
        =\prod_{n=1}^m \frac{n+1}{n}.
    \]

    This telescopes:
    \[
        p_m=\frac21\cdot\frac32\cdots\frac{m+1}{m}=m+1.
    \]

    Hence
    \[
        p_m\to\infty,
    \]
    so the product diverges.

    If $a_n=\frac1{n^2}$, then
    \[
        p_m=\prod_{n=1}^m\left(1+\frac1{n^2}\right).
    \]

    The first few terms are
    \[
        2,\quad 2\cdot\frac54,\quad 2\cdot\frac54\cdot\frac{10}{9},\quad \ldots
    \]

    This suggests convergence.

    \textbf{(b).}
    Assume first that
    \[
        \sum_{n=1}^\infty a_n
    \]
    converges.

    Let
    \[
        s_m=\sum_{n=1}^m a_n.
    \]

    Since $a_n\geq0$, the sequence $(s_m)$ is bounded.

    Also
    \[
        1+a_n\leq 3^{a_n}
    \]
    for $a_n\geq0$.

    Thus
    \[
        p_m=\prod_{n=1}^m(1+a_n)\leq \prod_{n=1}^m 3^{a_n}=3^{s_m}.
    \]

    So $(p_m)$ is bounded above.

    Because each factor $1+a_n\geq1$, $(p_m)$ is increasing.

    Hence $(p_m)$ converges.

    Conversely, assume $(p_m)$ converges.

    Then $(p_m)$ is bounded.

    We claim
    \[
        p_m\geq 1+\sum_{n=1}^m a_n.
    \]

    This is true by induction.

    For $m=1$,
    \[
        p_1=1+a_1.
    \]

    If it holds for $m$, then
    \begin{align*}
        p_{m+1}
        &=p_m(1+a_{m+1})\\
        &\geq \left(1+\sum_{n=1}^m a_n\right)(1+a_{m+1})\\
        &\geq 1+\sum_{n=1}^{m+1}a_n.
    \end{align*}

    Thus
    \[
        1+\sum_{n=1}^m a_n\leq p_m.
    \]

    Since $(p_m)$ is bounded, the partial sums of $\sum a_n$ are bounded.

    Because $a_n\geq0$, the partial sums are increasing.

    Hence
    \[
        \sum_{n=1}^\infty a_n
    \]
    converges.
\end{problem}

\begin{problem}[Infinite Products Two]{
    \begin{enumerate}[label={(\alph*)}]
\item Does
\[
    \frac21\cdot\frac32\cdot\frac54\cdot\frac98\cdot\frac{17}{16}\cdots
\]
converge?
\item Does
\[
    \frac12\cdot\frac34\cdot\frac56\cdot\frac78\cdot\frac9{10}\cdots
\]
converge to zero?
\item Show that
\[
    \left(\frac{2\cdot2}{1\cdot3}\right)\left(\frac{4\cdot4}{3\cdot5}\right)\left(\frac{6\cdot6}{5\cdot7}\right)\cdots
\]
at least converges.
    \end{enumerate}
}
    \textbf{(a).}
    This product is
    \[
        \prod_{n=1}^\infty\left(1+\frac{1}{2^n}\right).
    \]

    Using the identity
    \[
        \prod_{n=1}^m\left(1+\frac{1}{2^n}\right)=\frac{2-2^{-m}}{1},
    \]
    or by the standard telescoping trick,
    \[
        p_m=2-\frac{1}{2^m}.
    \]

    Hence
    \[
        p_m\to2.
    \]

    So the product converges.

    \textbf{(b).}
    This product is
    \[
        \prod_{k=1}^\infty \frac{2k-1}{2k}
        =
        \prod_{k=1}^\infty\left(1-\frac{1}{2k}\right).
    \]

    Each factor lies in $(0,1)$, so the partial products are decreasing and bounded below by $0$.

    Hence they converge.

    To determine the limit, note that
    \[
        \frac{1}{\prod_{k=1}^n \frac{2k-1}{2k}}
        =
        \prod_{k=1}^n \left(1+\frac{1}{2k-1}\right).
    \]

    But
    \[
        \sum_{k=1}^\infty \frac{1}{2k-1}
    \]
    diverges, so by the previous problem the reciprocal product diverges to $\infty$.

    Hence the original product converges to $0$.

    \textbf{(c).}
    The Wallis product can be written as
    \[
        \prod_{n=1}^\infty \frac{4n^2}{4n^2-1}
        =
        \prod_{n=1}^\infty \left(1+\frac{1}{4n^2-1}\right).
    \]

    Since
    \[
        0<\frac{1}{4n^2-1}\leq \frac{1}{n^2},
    \]
    and
    \[
        \sum_{n=1}^\infty \frac{1}{n^2}
    \]
    converges, it follows by comparison that
    \[
        \sum_{n=1}^\infty \frac{1}{4n^2-1}
    \]
    converges.

    Therefore the infinite product converges.
\end{problem}

\begin{problem}[Iterated Series]{
    Show that if
    \[
        \sum_{i=1}^{\infty}\sum_{j=1}^{\infty}\abs{a_{ij}}
    \]
    converges, then
    \[
        \sum_{i=1}^{\infty}\sum_{j=1}^{\infty} a_{ij}
    \]
    converges.
}
    For each fixed $i$, the series
    \[
        \sum_{j=1}^\infty \abs{a_{ij}}
    \]
    converges.

    Hence
    \[
        \sum_{j=1}^\infty a_{ij}
    \]
    converges absolutely, and therefore converges.

    Let
    \[
        b_i\coloneqq \sum_{j=1}^\infty \abs{a_{ij}}
        \qquad\text{and}\qquad
        c_i\coloneqq \sum_{j=1}^\infty a_{ij}.
    \]

    Then
    \[
        \abs{c_i}\leq b_i.
    \]

    Since
    \[
        \sum_{i=1}^\infty b_i
    \]
    converges, comparison implies
    \[
        \sum_{i=1}^\infty c_i
    \]
    converges.

    Thus
    \[
        \sum_{i=1}^{\infty}\sum_{j=1}^{\infty} a_{ij}
    \]
    converges.
\end{problem}